\section{A complete execution on the running instance}\label{sec:trace}

\subsection{Variable order and the eight pivots}
Order $y_1,\ldots,y_8$ first, followed by the nine slacks in the arc order
\eqref{eq:arcs}. Initialize with sink 1: $y_1=1$, all other vertex variables
zero, $\Fset=\emptyset$, and roots $2,\ldots,8$. The initial ratio is 1.
The exact reference execution is shown in \cref{tab:trace}.

\begin{table}[htbp]
\renewcommand{\arraystretch}{1.4}
\centering\input{../build/tutorial/generated/trace_table.tex}
\caption{Complete Bland-rule trace for the first ratio oracle. Each row is
checked against a separately assembled rational basis system.}\label{tab:trace}
\end{table}

Only pivots 4 and 7 have a positive step. The first three pivots make
dependencies tight before vertices 5 and 6 can increase. Pivots 5 and 6
also reorganize a zero region. Pivot 8 is needed to remove a positive
reduced cost from a degenerate optimal basis.

\subsection{A first zero step: a prerequisite blocks motion}
At the initial basis, $y_5$ has reduced cost $6-1\cdot2=4$.
Its raw displacement is $h=\ind^{\{5\}}$, and normalization gives
$d_5=1$, $d_1=-2$, with all other vertex directions zero. Consequently
\[
 d_{s_{5,1}}=-3,\qquad d_{s_{5,2}}=-1,\qquad d_{s_{5,6}}=-1.
\]
Although $s_{5,1}=1$, both $s_{5,2}$ and $s_{5,6}$ are zero. The step is
zero, and the fixed variable order makes $s_{5,2}$ leave. The new forest
contains $(5,2)$, its zero tree $\{2,5\}$ is rooted at 2, and the primal
point is unchanged. A positive reduced cost alone was insufficient to add
vertex 5: its prerequisites had to be handled first.

\subsection{The principal split: pivot 7}
Immediately before pivot 7, the positive tree and zero tree are
\[
 T_*=\{1,2,3,5,6\},\qquad T_0=\{4,7,8\},
\]
with $w(T_*)=6$, $p(T_*)=10$, and $\rho=5/3$. Their forest edges are
\[
 \Fset=\{(5,1),(5,2),(5,6),(6,3),(7,4),(8,7)\},
 \qquad \Qset=\{4\}.
\]
Root $T_*$ at 1. For $e=(5,6)$, parent 5 points to child 6, so
$\sigma_e=+1$ and $D_e=\{3,6\}$. Its aggregates are $p(D_e)=4$,
$w(D_e)=2$. Therefore
\[
 \bar c_{s_{5,6}}=4-\frac53\,2=\frac23>0,\qquad
 \kappa=\frac{2}{6}=\frac13.
\]
The vertex direction, in vertex order, is
\begin{equation}\label{eq:example-direction}
 d_y=\left(-\tfrac13,-\tfrac13,\tfrac23,0,-\tfrac13,\tfrac23,0,0\right).
\end{equation}
Each positive vertex initially has value $1/6$. The decreasing basic vertex
variables $y_1,y_2,y_5$ all give ratio $1/2$. So does the basic slack
$s_{8,5}=1/6$, whose direction is $d_5-d_8=-1/3$.
The smallest-index tied blocker is $y_1$, hence $t^\star=1/2$.
The update removes $(5,6)$ from $\Fset$ and adds root 1.

The new positive tree is $\{3,6\}$, with value $1/2$ at both vertices.
The vertices $\{1,2,5\}$ form a zero tree rooted at 1. The ratio increases to
\[
 \frac53+\frac12\cdot\frac23=2.
\]
\begin{figure}[htbp]
\centering
\begin{minipage}{.49\textwidth}\centering
\resizebox{\linewidth}{!}{\input{../build/tutorial/generated/forest_before.tex}}
\small Before pivot 7: positive tree $\{1,2,3,5,6\}$.
\end{minipage}\hfill
\begin{minipage}{.49\textwidth}\centering
\resizebox{\linewidth}{!}{\input{../build/tutorial/generated/forest_after.tex}}
\small After pivot 7: positive tree $\{3,6\}$.
\end{minipage}
\caption{The nonbasic-slack forest in solid blue, other arcs dashed.
Positive-tree vertices are shaded blue; double borders mark roots. The
right-hand basis still requires the degenerate pivot 8.}\label{fig:split}
\end{figure}

\subsection{An optimal point need not yet have an optimal basis}
After pivot 7, consider edge $(5,1)$ in the zero tree $\{1,2,5\}$,
rooted at 1. It points from child 5 to parent 1, so $\sigma=-1$ and
$D=\{2,5\}$. Pricing gives
\[
 \bar c_{s_{5,1}}=-\bigl[5-2\cdot3\bigr]=1>0.
\]
Its raw motion decreases $y_2$ and $y_5$, which are already zero. The
ratio test returns zero and $y_2$ leaves. The objective remains 2, but
the forest now has roots $1,2,4$ and edges
\[
 \Fset=\{(5,2),(6,3),(7,4),(8,7)\}.
\]
All nonbasic reduced costs are now nonpositive. The ratio certificate is
\begin{equation}\label{eq:ratio-example-cert}
 \beta=2,\qquad
 \alpha_{5,2}=2,\quad\alpha_{6,3}=3,\quad
 \alpha_{7,4}=1,\quad\alpha_{8,7}=1,
\end{equation}
with every other arc multiplier zero. Its vertex residuals are 1, 1 and 4
at roots 1, 2 and 4, respectively, and zero at all other vertices.
It proves that no nonempty closure has ratio greater than 2.

This final pivot is a useful review case: stopping solely because the
objective did not change, or because the support already looks correct,
would bypass the basis optimality test.
